% Auto-generated by src/analysis/latex_tables.py — do not edit by hand.
\begin{table}[htbp]
\centering
\small
\caption{Variance decomposition: $R^{2}$ of nested specifications and marginal contribution of the treatment}
\label{tab:variance_decomposition}
\begin{tabular}{lccccc}
\toprule
 & \multicolumn{4}{c}{$R^{2}$ of nested specifications} & \multicolumn{1}{c}{Marginal} \\
\cmidrule(lr){2-5}\cmidrule(lr){6-6}
Dependent variable & County FE & Year FE & Both FE & $+$ Treatment & contribution \\
\midrule
Crude birth rate & 0.658 & 0.164 & 0.803 & 0.804 & 0.0008 \\
Legit.\ birth rate & 0.682 & 0.139 & 0.806 & 0.807 & 0.0012 \\
Illegitimacy ratio & 0.906 & 0.013 & 0.917 & 0.918 & 0.0006 \\
Gen.\ marriage rate & 0.345 & 0.195 & 0.537 & 0.540 & 0.0030 \\
\bottomrule
\end{tabular}
\begin{minipage}{\linewidth}
\vspace{0.5em}
\footnotesize \textit{Notes:} Each cell reports the $R^{2}$ of the named specification. ``Marginal contribution'' is the gain in $R^{2}$ from adding $\mathrm{CathShare} \times \mathrm{Post}$ to the two-way FE specification. County FE absorb the bulk of the variation in fertility outcomes; the marginal contribution of the treatment is small in absolute terms but informative about identification: a near-zero marginal $R^{2}$ alongside a precisely estimated coefficient (as in the marriage-rate column) implies a real but narrow within-county effect.
\end{minipage}
\end{table}
